\subsection{Variational structure, spectral assignment, and local curvature}
Each prefix of the flag is evaluated against a covariance with its own memory
scale. Consequently a direction assigned to an inner block contributes to
several prefix objectives, as encoded by $A_j$. The assignment theorem below
resolves these competing contributions under simultaneous diagonalization,
while respecting the prescribed capacity of every block.

\begin{definition}[Flag energy]
\begin{gather}
 C_k=C_k^\top\succeq0,\quad \omega_k>0,\quad
 A_j:=\sum_{k=j}^K\omega_k C_k,\quad A_{K+1}:=0,\\
 \Phi(U;C):=\sum_{k=1}^K\omega_k\tr(P_k(U)C_k)
 =\sum_{j=1}^K\tr(U_j^\top A_jU_j),
 \label{eq:tf-objective}\\
 H(U;C):=2\sum_k\omega_k C_kUF_k,\quad
 \operatorname{grad}\Phi(U;C)=\mathsf H_U(H(U;C)).
\end{gather}
\end{definition}
\begin{proposition}[First and second variations]\label{prop:rigid-variation}
For $Z\in\Hor_U$, $Q=[U,U_\perp]$, and horizontal $\Omega$ with
$Z=(Q\Omega)_{:R}$,
\begin{align}
 \mathrm D\Phi(U)[Z]&=2\sum_k\omega_k\tr(Z^\top C_kUF_k)
 =\langle\operatorname{grad}\Phi,Z\rangle_F,\\
 \left.\frac{\mathrm d^2}{\mathrm ds^2}\Phi((Qe^{s\Omega})_{:R};C)\right|_{s=0}
 &=\sum_k\omega_k\tr\bigl(Q^\top C_kQ[\Omega,[\Omega,\widehat F_k]]\bigr),\\
 \operatorname{grad}\Phi(U;C)=0
 &\Longrightarrow\quad
 \langle Z,\operatorname{Hess}\Phi(U)[Z]\rangle_F
 =\sum_k\omega_k\tr\bigl(Q^\top C_kQ[\Omega,[\Omega,\widehat F_k]]\bigr).
\end{align}
\end{proposition}
\begin{proof}
\begin{align}
 \mathrm D(UF_kU^\top)[Z]&=ZF_kU^\top+UF_kZ^\top,\\
 \frac{\mathrm d}{\mathrm ds}(e^{s\Omega}\widehat F_ke^{-s\Omega})
 &=e^{s\Omega}[\Omega,\widehat F_k]e^{-s\Omega},\\
 \frac{\mathrm d^2}{\mathrm ds^2}(e^{s\Omega}\widehat F_ke^{-s\Omega})
 &=e^{s\Omega}[\Omega,[\Omega,\widehat F_k]]e^{-s\Omega},\\
 (\Phi\circ\gamma)''(0)
 &=\operatorname{Hess}\Phi[\dot\gamma(0),\dot\gamma(0)]
   +\langle\operatorname{grad}\Phi,\nabla_{\dot\gamma}\dot\gamma\rangle_F.
\end{align}
\end{proof}
\begin{theorem}[Commuting case: integral assignment]\label{thm:tf-assignment}
Suppose
\begin{gather}
 C_k=Q\diag(\lambda_{k1},\ldots,\lambda_{kd})Q^\top,\quad
 b_{ji}:=\sum_{k=j}^K\omega_k\lambda_{ki},\quad b_{K+1,i}:=0,\\
 \mathcal X:=\{x\ge0:\sum_{j=1}^{K+1}x_{ji}=1,
                    \ \sum_{i=1}^d x_{ji}=q_j\}.
\end{gather}
Then
\begin{align}
 \max_{[U]\in\mathfrak F}\Phi(U;C)
 &=\max_{x\in\mathcal X}\sum_{j,i}b_{ji}x_{ji}\\
 &=\max_{\substack{\{S_j\}\text{ partition }[d]\\ |S_j|=q_j}}
       \sum_{j=1}^{K+1}\sum_{i\in S_j}b_{ji}\\
 &=\min_{\alpha\in\RR^d,\,\beta\in\RR^{K+1}}
       \left\{\sum_i\alpha_i+\sum_jq_j\beta_j:
                    \alpha_i+\beta_j\ge b_{ji}\right\}.
\end{align}
\end{theorem}
\begin{proof}
\begin{align}
 x_{ji}(U)&:=e_i^\top Q^\top\Pi_j(U)Qe_i\ge0,\\
 \sum_jx_{ji}(U)&=1,\qquad \sum_ix_{ji}(U)=\tr\Pi_j=q_j,\\
 \Phi(U;C)&=\sum_{j,i}b_{ji}x_{ji}(U)
 \le\max_{x\in\mathcal X}\langle b,x\rangle.
\end{align}
Changing the signs of the block-sum rows turns the constraint matrix into a
directed bipartite incidence matrix. It is totally unimodular, so the
Hoffman--Kruskal integrality theorem~\citep{hoffman1956integral}, with integer
margins $(q_j,1)$, gives integral vertices. Column sums of one then give
$x_{ji}\in\{0,1\}$. For $S_j:=\{i:x_{ji}=1\}$ choose
$U_j=Q_{:,S_j}$. Then $x(U)=x$. Finally,
\begin{align}
 \mathcal L(x,\alpha,\beta)
 &=\sum_i\alpha_i+\sum_jq_j\beta_j+
       \sum_{j,i}(b_{ji}-\alpha_i-\beta_j)x_{ji},\\
 \sup_{x\ge0}\mathcal L<\infty
 &\Longleftrightarrow\alpha_i+\beta_j\ge b_{ji},\\
 \varnothing\ne\mathcal X\text{ compact}
 &\Longrightarrow\max\langle b,x\rangle=\min\mathcal L_{\mathrm{dual}}.
\end{align}
\end{proof}
\begin{corollary}[Common order and strict assignment margin]\label{cor:tf-common-order}
If $\lambda_{k1}\ge\cdots\ge\lambda_{kd}$ for every $k$, then
$U^\star=Q_{:,1:R}$ maximizes $\Phi$. More generally, let an optimal partition
$S^\star$ admit a dual optimum satisfying
\begin{gather}
 s_{ji}:=\alpha_i+\beta_j-b_{ji},\quad
 s_{ji}=0\ (i\in S_j^\star),\quad s_{ji}\ge\delta>0\ (i\notin S_j^\star),\\
 W(U):=\sum_j\sum_{i\notin S_j^\star}x_{ji}(U).
\end{gather}
Then
\begin{align}
 \Phi(U^\star)-\Phi(U)&=\sum_{j,i}s_{ji}x_{ji}(U)\ge\delta W(U),\\
 d_P(U,U^\star)^2&\le K W(U),\\
 \Phi(U^\star)-\Phi(U)&\ge\frac{\delta}{K}d_P(U,U^\star)^2.
\end{align}
The factor $K$ in the distance comparison is attained by a rotation between
the first and residual blocks.
\end{corollary}
\begin{proof}
Ky Fan's maximum principle~\citep{fan1949weyl} gives the common-order optimum.
For the quantitative bound let $T_{ja}:=\tr(\Pi_j(U)\Pi_a(U^\star))$, including
the residual block. Its row and column sums both equal the block capacities.
\begin{align}
 \tr(P_kC_k)&\le\sum_{i=1}^{r_k}\lambda_{ki}
 =\tr(P_k(U^\star)C_k),\\
 \|\Pi_j(U)-\Pi_j(U^\star)\|_F^2
 &=2q_j-2\sum_{i\in S_j^\star}x_{ji}
 =2\sum_{i\notin S_j^\star}x_{ji},\\
 \|P_k(U)-P_k(U^\star)\|_F^2
 &=2\sum_{j\le k<a}T_{ja}
 =\sum_{j\le k<a}T_{ja}+\sum_{a\le k<j}T_{ja},\\
 d_P^2&=\sum_{j,a}|j-a|T_{ja}\le K\sum_{j\ne a}T_{ja}=KW(U).
\end{align}
For a rotation of angle $\theta$ between a vector of block $1$ and one of
block $K+1$, $W=2\sin^2\theta$ and $d_P^2=2K\sin^2\theta$. Thus the geometric
constant cannot be reduced. The dual-slack identity gives the objective bound.
\end{proof}
\begin{proposition}[Perturbation and matched-rank reconstruction]
\label{prop:tf-value}
For $\|\widehat C_k-C_k\|_{\op}\le\epsilon_k$, define
$D_C:=\sum_k\omega_kr_k\epsilon_k$. If $U^\star,\widehat U$ maximize their
respective objectives, then
\begin{gather}
 \sup_U|\Phi(U;\widehat C)-\Phi(U;C)|\le D_C,\qquad
 0\le\Phi(U^\star;C)-\Phi(\widehat U;C)\le2D_C.
\end{gather}
For every rank-$R$ projector $Q_R$ containing $\operatorname{range}P_k$,
\begin{equation}\label{prop:tf-budget}
 \|G-GQ_R\|_F^2\le\|G-GP_k\|_F^2.
\end{equation}
\end{proposition}
\begin{proof}
\begin{align}
 |\tr(P_k\Delta C_k)|&\le\|P_k\|_*\|\Delta C_k\|_{\op}=r_k\epsilon_k,\\
 \Phi(U^\star;C)-\Phi(\widehat U;C)
 &\le\Phi(U^\star;C)-\Phi(U^\star;\widehat C)
       +\Phi(\widehat U;\widehat C)-\Phi(\widehat U;C),\\
 Q_R-P_k&\succeq0,\quad
 \|G-GP_k\|_F^2-\|G-GQ_R\|_F^2
 =\tr((Q_R-P_k)G^\top G)\ge0.
\end{align}
\end{proof}
\begin{proposition}[A strict nonglobal local maximum]\label{prop:tf-cycle}
Let
\begin{gather}
 d=3,\quad K=2,\quad(r_1,r_2)=(1,2),\quad\omega_1=\omega_2=1,\\
 C_1=\diag(3,\tfrac52,\tfrac12),\quad
 C_2=\diag(0,2,\tfrac32),\quad
 (b_{ji})=\begin{pmatrix}3&9/2&2\\0&2&3/2\\0&0&0\end{pmatrix},\\
 U_0=(e_1,e_2),\qquad U_\star=(e_2,e_3).
\end{gather}
Then $[C_1,C_2]=0$, $\Phi(U_0)=5<6=\Phi(U_\star)$, and
\begin{equation}
 \Phi((e^{\Omega})_{:2})=5-\tfrac12\Omega_{12}^2
                              -\Omega_{13}^2-\tfrac12\Omega_{23}^2
                              +O(\|\Omega\|_F^3).
\end{equation}
\end{proposition}
\begin{proof}
For eigen-indices $i\in S_a$, $l\in S_b$, $a<b$, put
\begin{gather}
 \Delta_{ai,bl}:=b_{ai}+b_{bl}-b_{al}-b_{bi},\\
 \left.\frac{\mathrm d^2}{\mathrm ds^2}\Phi((Qe^{s\Omega})_{:R})\right|_{0}
 =-2\sum_{a<b}\sum_{i\in S_a,l\in S_b}\Delta_{ai,bl}\Omega_{il}^2.
\end{gather}
At $U_0$,
\begin{gather}
 \Delta_{11,22}=\tfrac12,\quad \Delta_{11,33}=1,\quad
 \Delta_{22,33}=\tfrac12,\quad
 \operatorname{grad}\Phi(U_0)=0,\quad\operatorname{Hess}\Phi(U_0)\prec0.
\end{gather}
The six assignment scores are $5,9/2,9/2,6,2,4$.
\end{proof}
The counterexample distinguishes global allocation from local improvement:
every infinitesimal pairwise rotation at $U_0$ decreases the objective, although
a different assignment has greater value. Commutativity therefore provides an
exact global characterization without ensuring that ascent finds it. The
oracle used subsequently is an interior stationary branch with uniform local
curvature; the analysis requires the iterates to remain in its regular chart.

\begin{assumption}[A regular oracle chart]\label{ass:tf-chart}
Let $\psi:\mathcal D\subset\RR^p\to\mathfrak F$ be a smooth chart on a convex
set, $f_C:=\Phi(\psi(\cdot);C)$, and
\begin{gather}
 \nabla f_C(x_C^\star)=0,\quad
 -L_\phi I\preceq\nabla^2 f_C(x)\preceq-\mu_\phi I\quad(x\in\mathcal D),\\
 0<\mu_\phi\le L_\phi,\quad
 \|\mathrm D\psi(x)\|_{2\to F}\le J_\psi,\quad
 \|\nabla^3 f_C(x)\|\le T_\phi.
\end{gather}
\end{assumption}
\begin{lemma}[Oracle sensitivity and differentiability]\label{lem:tf-oracle}
If $C,\widehat C$ satisfy Assumption~\ref{ass:tf-chart} in the same chart,
then, with $E_C:=2J_\psi\sum_k\omega_k\sqrt{r_k}\epsilon_k$,
\begin{align}
 \sup_x\|\nabla f_{\widehat C}(x)-\nabla f_C(x)\|&\le E_C,\\
 \|x_{\widehat C}^\star-x_C^\star\|&\le E_C/\mu_\phi,\\
 \mathrm D x_C^\star[\Delta C]
 &=-[\nabla^2 f_C(x_C^\star)]^{-1}
                 \nabla f_{\Delta C}(x_C^\star).
\end{align}
If $\|\nabla^2f_{\Delta C}\|\le E_H$, and $v=x_{C+\Delta C}^\star-x_C^\star$,
\begin{equation}
 \|v-\mathrm D x_C^\star[\Delta C]\|
 \le\frac{T_\phi}{2\mu_\phi^3}E_C^2+
                         \frac{E_H}{\mu_\phi^2}E_C.
\end{equation}
\end{lemma}
\begin{proof}
\begin{align}
 \|H(U;\Delta C)\|_F&\le2\sum_k\omega_k\epsilon_k\|UF_k\|_F
 =2\sum_k\omega_k\sqrt{r_k}\epsilon_k,\\
 \mu_\phi\|v\|^2
 &\le-\langle\nabla f_C(x_C^\star+v)-\nabla f_C(x_C^\star),v\rangle
 \le E_C\|v\|,\\
 0&=\nabla f_{C+\Delta C}(x_C^\star+v)
 =\nabla^2 f_C(x_C^\star)v+\nabla f_{\Delta C}(x_C^\star)+R,\\
 \|R\|&\le\tfrac12T_\phi\|v\|^2+E_H\|v\|,\\
 v&=-[\nabla^2 f_C(x_C^\star)]^{-1}
                  (\nabla f_{\Delta C}(x_C^\star)+R).
\end{align}
\end{proof}
The factor $\mu_\phi^{-1}$ measures how strongly estimation errors can move the
local flag optimum. A nearly flat objective can change its preferred flag
substantially under a small covariance perturbation. The next spectral bound
exhibits the corresponding gap dependence for a single projector; a general
multi-covariance flag uses the chart curvature above.

\begin{proposition}[Residual Davis--Kahan bound]\label{prop:rigid-gap}
Let $C=Q_1\Lambda_1Q_1^\top+Q_2\Lambda_2Q_2^\top$, $\dim Q_1=r$, and
$g:=\lambda_r(C)-\lambda_{r+1}(C)>0$. If $\widehat C=C+E$,
$\|E\|_{\op}\le\epsilon<g$, and $\widehat P$ is its leading rank-$r$ projector,
set $s_r:=\min(r,d-r)$ and
$g_{\rm sep}:=\lambda_{\min}(\widehat\Lambda_1)-\lambda_{\max}(\Lambda_2)$.
Then $g_{\rm sep}\ge g-\epsilon>0$ and
\begin{equation}
 \|\widehat P-Q_1Q_1^\top\|_F
 \le\frac{\sqrt2\|Q_2^\top E\widehat Q_1\|_F}{g_{\rm sep}}
 \le\frac{\sqrt{2s_r}\,\epsilon}{g-\epsilon}.
\end{equation}
The left side is also at most $\sqrt{2s_r}$. Keeping the off-diagonal residual
can be substantially sharper than replacing it by $\sqrt{s_r}\epsilon$.
\end{proposition}
\begin{proof}
This is the Sylvester-equation form of the Davis--Kahan
$\sin\Theta$ theorem~\citep{davis_kahan}; Weyl's inequality supplies the
separation from the unperturbed gap.
\begin{align}
 \widehat C\widehat Q_1&=\widehat Q_1\widehat\Lambda_1,\quad
 X:=Q_2^\top\widehat Q_1,\quad
 X\widehat\Lambda_1-\Lambda_2X=Q_2^\top E\widehat Q_1=:B,\\
 \lambda_{\min}(\widehat\Lambda_1)-\lambda_{\max}(\Lambda_2)
 &\ge g-\epsilon,\\
 X_{ij}&=B_{ij}/(\widehat\lambda_j-\lambda_{r+i}),\quad
 \|X\|_F\le\|B\|_F/g_{\rm sep},\quad
 \|B\|_F\le\sqrt{s_r}\|E\|_{\op},\\
 \|\widehat P-P\|_F^2&=2r-2\tr(\widehat P P)=2\|Q_2^\top\widehat Q_1\|_F^2.
\end{align}
\end{proof}
\begin{corollary}[Local curvature at a common-order flag]
If $\lambda_{k1}\ge\cdots\ge\lambda_{kd}$, then for $i\in I_a$, $l\in I_b$,
\begin{gather}
 \Delta_{ai,bl}=\sum_{k=a}^{b-1}\omega_k(\lambda_{ki}-\lambda_{kl}),
 \quad \lambda_{k\cdot}:=0\ (k>K),\\
 \delta_\star:=\min_{a<b,\,i\in I_a,\,l\in I_b}\Delta_{ai,bl}>0
 \ \Longrightarrow\quad
 -\operatorname{Hess}\Phi(U^\star)[Z,Z]\ge\delta_\star\|Z\|_F^2.
\end{gather}
For orthonormal chart coordinates at $U^\star$, and
$\|\nabla^2f(x)-\nabla^2f(0)\|\le T_\phi\|x\|$,
\begin{equation}
 \|x\|\le\delta_\star/(2T_\phi)
 \quad\Longrightarrow\quad-\nabla^2f(x)\succeq(\delta_\star/2)I.
\end{equation}
\end{corollary}
